% storage accounting (static text; numbers via \val)
Storage is the logical state a hardware implementation would hold, in bits. For the TAGE family the predictor counts its own state (\code{storage\_bits()}) and prints it at start-up; the printed value was checked in every one of the corresponding runs. The ChampSim predictors print nothing, so their state was counted from the source code: tables and history registers are counted at their declared widths, the 8-bit width of the hashed perceptron's adaptive threshold is an assumption, and the perceptron's speculative copy of the history is not counted because it is identical to the architectural copy under the immediate update of this ChampSim. Table~\ref{tab:breakdown} gives the breakdown of \prop{}.

Three storage classes result: about 4 KB (bimodal, gshare, perceptron), 47 KB (TAGE) and 58--64 KB (TAGE-SC-L, \prop{}, hashed perceptron). \prop{} uses \val{prop_bits} bits (\val{prop_kb} KB): \val{tsl_kb}~KB for the TAGE-SC-L part plus \val{prop_minus_tsl_bits} bits (\val{prop_minus_tsl_kb} KB, +\val{prop_minus_tsl_pct}\%) for the loop-structure components and the inference tables. It is \val{hp_minus_prop_bits} bits smaller than the 64 KB hashed perceptron (\val{hp_bits} bits). The like-for-like comparisons are therefore with the hashed perceptron (largest ChampSim predictor, similar budget) and TAGE-SC-L (same implementation without the contribution); the 4 KB predictors are included because they are ChampSim's defaults, not as storage-equivalent competitors.

Storage in bits is not silicon area. ChampSim models neither predictor access latency nor area or energy, so these were not measured. Qualitatively, \prop{} reads 13 TAGE tables, a loop table and 18 corrector tables per prediction (the hashed perceptron reads 16 tables); TAGE-SC-L designs are normally pipelined over two cycles with an overriding scheme, and the extra corrector inputs lengthen the adder tree by about one level.
